\documentclass[12pt]{article}

%%% These are some packages that are useful
\usepackage{amsfonts, lipsum}
\usepackage[scr]{rsfso}
\usepackage{amsmath,amssymb, amscd,amsbsy, bbm, amsthm, enumerate}
\usepackage[top=1in, bottom=1in, left=.45in, right=.45in]{geometry}
\usepackage[unicode]{hyperref}
\usepackage{tikz, xcolor, fancyhdr}
\usepackage{graphicx, caption, subcaption}

%%% Page formatting
%\setlength{\headsep}{30pt}
\setlength{\parindent}{0pt}
\setlength{\textheight}{9in}



%%% These define our question environment and help number things correctly
\theoremstyle{definition}
\newtheorem{thm}{Theorem}
\newtheorem{question}[thm]{Question}
\newtheorem{questionS}[thm]{***Question}
\newtheorem{prop}[thm]{Proposition}
\newtheorem{lem}[thm]{Lemma}
\newtheorem{cor}[thm]{Corollary}
\newtheorem{defn}[thm]{Definition}
\newtheorem{rem}[thm]{Remark}
\newtheorem{ex}[thm]{Example}

%%% These are some shortcuts that are handy
\def\real{{\mathbb R}}
\def\Natural{\mathbb{N}}
\def\dx{\textnormal{dx}}
\def\dy{\textnormal{dy}}
\def\dz{\textnormal{dz}}
\def\dt{\textnormal{dt}}
\def\Re{\textnormal{Re}}
\def\Im{\textnormal{Im}}
\def\exp{\textnormal{exp}}
\def\interior{\textnormal{interior}}
\def\al{\alpha}
\def\del{\delta}
\def\Del{\Delta}
\def\gam{\gamma}
\def\Gam{\Gamma}
\def\Om{\Omega}
\def\ep{\varepsilon}
\def\lam{\lambda}
\def\rational{{\mathbb Q}}
\def\integer{{\mathbb Z}}
\def\Q{{\mathbb Q}}
\def\Z{{\mathbb Z}}
\def\N{{\mathbb N}}
\def\R{{\mathbb R}}
\def\grad{\nabla}
\def\C{\mathcal C}
\def\P{\mathbb{P}}
\def\T{\mathcal T}
\def\I{\mathcal I}
\newcommand{\abs}[1]{\left| #1 \right|}
\newcommand{\inner}[1]{\langle #1 \rangle}
\newcommand{\norm}[1]{\left\lVert#1\right\rVert}
\newcommand{\spanvect}{\textnormal{span}}
\newcommand{\union}{\cup}
\newcommand{\Union}{\bigcup}
\def\intersect{\cap}
\def\Intersect{\bigcap}
\renewcommand{\neg}{\mathord{\sim}}
\newcommand{\ds}{\displaystyle}





%%%%%%%%%%%%%%%%%%%%%%%%%%%%%%%%%
%%%%%%%%%%%%%%%%%%%%%%%%%%%%%%%%%
\newcommand{\name}{Duff Baker-Jarvis}
%%%%%%%%%%%%%%%%%%%%%%%%%%%%%%%%%
%%%%%%%%%%%%%%%%%%%%%%%%%%%%%%%%%
%%% Document Starts now
\begin{document}
\title{Analysis - Lay (OUTLINE)}
\date{\today}
\author{Dr. Duff Baker-Jarvis}
\maketitle


\begin{enumerate}
    \item Section 13 - Topology of the Reals
	\begin{itemize}
	    \item (Defn) Neighborhood and Deleted Neighborhood
	    \item (Defn) Interior point
	    \item (Defn) Boundary point
	    \item (Defn) Open and Closed Sets
	    \item (Thm) Equivalent definitions of open and closed sets
	    \item (Thm) Arbitrary unions of open sets are open\\
		    Proof idea: Every point in the union lies in some open set. Take a neighborhood that lies within that set. It also lies within the union.
	    \item (Thm) Finite intersections of open sets are open\\
		Proof idea: Let $ c $ be in the finite intersection. Then for each set there is a neighborhood of $ c $ contained within that set. Each has a radius. 
		consider the neighborhood with the smallest raidus. Then  this lies in the intersection.
	    \item (Thm) Arbitray intersections of closed sets are closed.\\
		Proof idea: Way 1 - by complements of previous theorem. Use Demorgan's law. Way 2 - by defininiton of boundary.
	    \item (Thm) Finite unions of closed sets are closed.
	\end{itemize}
	
\end{enumerate}



\end{document}
